% Venue-neutral draft. To switch to a venue template, replace the
% \documentclass + layout block below; content lives in sections/ and
% macros.tex, which use only standard packages.
\documentclass[10pt]{article}
\usepackage[margin=1in]{geometry}

\usepackage[T1]{fontenc}
\usepackage{lmodern}
\usepackage{amsmath,amssymb}
\usepackage{booktabs}
\usepackage{graphicx}
\usepackage{xcolor}
\usepackage[numbers,sort&compress]{natbib}
\usepackage[hidelinks]{hyperref}

\newcommand{\todo}[1]{{\color{red}\textbf{[TODO:} #1\textbf{]}}}
% stands in for a table or figure that is not produced yet
\newcommand{\placeholder}[1]{%
  \noindent\fbox{\parbox{0.95\linewidth}{\centering\small\color{gray}%
  \textbf{PLACEHOLDER}\\#1}}}
% tables/<name>.tex is a bare tabular written by the results generator
\newcommand{\resulttable}[2]{%
  \IfFileExists{tables/#1.tex}{\input{tables/#1.tex}}{\placeholder{#2}}}

\newcommand{\plr}{PLR$^\perp$}
\newcommand{\cvar}[1]{CVaR$_{#1}$}


\title{Unsupervised Environment Design for Legged Locomotion:\\
A Minimal Benchmark on Closed- and Open-Chain Quadrupeds}
\author{Anonymous Authors}
\date{}

\begin{document}
\maketitle

\begin{abstract}
Unsupervised environment design (UED) methods such as PLR, ACCEL, SFL and
PAIRED are mostly compared on 2D grid worlds, mazes and BipedalWalker, while
legged-locomotion curricula are usually hand-designed and evaluated against
a single baseline. To our knowledge, no benchmark compares several UED
methods under one protocol on 3D legged locomotion.
We present a minimal one: two quadrupeds in MuJoCo MJX, a small closed-chain
robot (Wojtek, ${\sim}0.10$\,m tall) and the open-chain Unitree Go1, a
level space of eight interpretable terrain parameters, a brax PPO student
trained in a loop that swaps the heightfield of every environment at every
reset, and six methods (DR, DR-easy, PLR$^\perp$, ACCEL, ACCEL-SFL, PAIRED)
at matched environment-step budgets.
The evaluation protocol reports uniform-level success, cross-fitted CVaR,
per-level pass-to-pass reliability, per-feature capability envelopes and a
held-out structured terrain suite.
The level space is calibrated with pilot policies so that uniform sampling is
neither trivially easy nor hopeless.
In a pilot on a wider, uncalibrated space, the starting level distribution
dominated: DR on easy levels matched ACCEL, and regret-based curation
recovered part of the gap left by a uniform generator.
\todo{Replace the last sentence with the v10 headline once results land.}
\end{abstract}

\section{Introduction}
\label{sec:intro}

Learned locomotion controllers for legged robots are trained on
distributions of terrain, and the choice of that distribution shapes what
the controller can do. Practical pipelines use hand-designed
curricula: terrain types arranged by a difficulty parameter that a robot is
promoted along once it walks far enough~\citep{rudin2022learning}, or
adaptive samplers over a parameterised terrain
family~\citep{lee2020learning}. Recent locomotion work automates this with
learning-progress~\citep{li2026lpacrl}, learned-difficulty~\citep{liu2026tacl}
and predicted-safety-cost~\citep{ovinnikov2026psc} curricula, each evaluated
with its own baselines, robot and protocol.

Unsupervised environment design (UED) studies the same problem in general
form: a teacher proposes environment parameters (\emph{levels}) to a student
so that the student generalises, typically by maximising an estimate of the
student's regret. PAIRED~\citep{dennis2020paired} learns a level generator;
PLR and its robust variant \plr~\citep{jiang2021plr,jiang2021replay} curate
randomly sampled levels by a regret score; ACCEL~\citep{parkerholder2022accel}
adds mutation of curated levels; SFL~\citep{rutherford2024noregrets} replaces
regret estimates with learnability. These methods are compared mostly on
mazes, grid worlds, car racing and the 2D BipedalWalker; libraries such as
minimax~\citep{jiang2023minimax} and JaxUED~\citep{coward2024jaxued} make
those comparisons cheap. To our knowledge no benchmark runs several UED
methods under one protocol on 3D legged locomotion, where the student is a
high-dimensional, contact-rich controller and a level is a terrain.

We build a minimal one and use it to ask whether the UED findings from 2D
domains carry over. Our contributions are:
\begin{itemize}
  \item \textbf{A two-robot benchmark} in MuJoCo MJX~\citep{todorov2012mujoco}
    with a closed-kinematic-chain quadruped (Wojtek, ${\sim}0.10$\,m
    standing height, our own physically built robot) and the open-chain
    Unitree Go1 task from MuJoCo Playground~\citep{zakka2025playground},
    scaled so that both face the same terrain relative to body size
    (\S\ref{sec:benchmark}).
  \item \textbf{An interpretable, calibrated level space}: eight terrain
    parameters plus a seed, whose ranges are set from the capability
    envelopes of pilot policies, aiming for uniform sampling that is
    neither hopeless nor trivially easy (\S\ref{sec:levelspace}).
  \item \textbf{Six methods on one student}: DR, DR-easy, \plr, ACCEL,
    ACCEL-SFL and PAIRED, all driving the same brax
    PPO~\citep{schulman2017ppo,freeman2021brax} student at matched
    environment-step budgets (\S\ref{sec:methods}).
  \item \textbf{An evaluation protocol} with cross-fitted CVaR, a
    cross-fitted solvable-subset success, per-level reliability,
    capability envelopes, a held-out structured terrain suite and a
    sim-to-sim transfer check (\S\ref{sec:protocol}), and a pilot study
    whose outcome motivated the calibration (\S\ref{sec:results},
    Appendix~\ref{app:pilot}).
\end{itemize}
A method interacts with the benchmark only through the level vectors it
proposes and whether the student trains on them, so the level space,
calibration and evaluation carry over unchanged to new methods.
\todo{Code release statement (anonymised link) and licence.}

\section{Related Work}
\label{sec:related}

\paragraph{Unsupervised environment design.}
PAIRED~\citep{dennis2020paired} frames curriculum design as a game between
a learned level generator and two students, a protagonist and an
antagonist, whose return gap approximates regret.
PLR~\citep{jiang2021plr} keeps a buffer of sampled levels prioritised by a
score such as positive value loss and replays them by rank and staleness;
\plr~\citep{jiang2021replay} trains only on replayed levels, which makes it
a regret-minimax method.
ACCEL~\citep{parkerholder2022accel} extends \plr\ by editing replayed
levels and scoring the edits as new candidates, with a generator of simple
levels, and reports gains on 2D BipedalWalker terrain.
SFL~\citep{rutherford2024noregrets} shows that common regret approximations
track success rate rather than regret, scores levels by learnability
$p(1-p)$ for success probability $p$, and proposes an adversarial
CVaR-style evaluation that we adapt.
minimax~\citep{jiang2023minimax} and JaxUED~\citep{coward2024jaxued}
provide accelerated reference implementations on 2D domains.
Our benchmark keeps these algorithms close to their published form and
changes the domain: levels are heightfields for a 3D quadruped.

\paragraph{Terrain curricula for legged locomotion.}
\citet{lee2020learning} train ANYmal over parameterised terrain with an
adaptive curriculum that samples terrain parameters of intermediate
difficulty. \citet{rudin2022learning} introduce a game-inspired curriculum
in which robots move to harder terrain rows after walking far enough; it is
the default in massively parallel locomotion training.
LP-ACRL~\citep{li2026lpacrl} samples tasks by estimated learning progress
and reaches high-speed ANYmal~D locomotion across stairs, slopes, gravel and
low-friction ground.
TACL~\citep{liu2026tacl} learns a difficulty function over
trajectory-level tasks on unstructured terrain maps and proposes the next
tasks from it.
PSC~\citep{ovinnikov2026psc} prioritises terrain contexts and randomised
events by a distributional safety critic and compares against terrain
progression, advantage-based replay and learning-progress curricula.
Eurekaverse~\citep{liang2024eurekaverse} uses an LLM to write progressively
harder parkour courses for a quadruped.
\citet{zhang2022terrain} generate terrain datasets by authoring and active
learning as a prototype robustness benchmark.
Each of these works proposes one curriculum or dataset and compares it
with a few baselines; none runs the standard UED family (PAIRED, \plr,
ACCEL, SFL) side by side under a shared level space and evaluation.
\todo{Final literature sweep before submission (2025--2026 UED and
locomotion-curriculum papers); check whether TACL or PSC include PLR/ACCEL
baselines in their experiments, not only in related work.}

\section{Benchmark}
\label{sec:benchmark}

\subsection{Robots and task}
\label{sec:robots}

\paragraph{Wojtek} is a small quadruped built by the authors, with a
standing height of about 0.10\,m and 12 actuated joints (three per leg).
Each leg is a closed kinematic chain: the actuated joints drive the foot
through a passive four-bar linkage, closed in MuJoCo by one
\texttt{connect} equality constraint per leg. The closure adds constraint
dynamics that serial legs do not have (the linkage can reach a toggle
configuration, which the reward penalises), which is why we pair Wojtek
with an open-chain robot.
\todo{Add mass, link lengths and a photo of the physical robot.}

\paragraph{Unitree Go1} is the open serial-chain quadruped of the MuJoCo
Playground joystick task~\citep{zakka2025playground}, with Playground's
reward, feet-only collisions and physics randomisation. Its standing trunk
height is 0.258\,m. We change four things so it runs the same benchmark:
the scene is our heightfield arena; commands are forward-walking;
termination uses Wojtek's fall rule at Go1's size (trunk below
$0.6\times0.258$\,m above the terrain under it, or gravity $z$-component in
the body frame above $-0.4$); and the critic receives a terrain height scan.

\paragraph{Task.}
Both robots track a velocity command
$(v_x, v_y, \omega_z)$ drawn per episode from
$v_x\in[0.2,0.6]$, $v_y\in[-0.2,0.2]$\,m/s and
$\omega_z\in[-0.4,0.4]$\,rad/s for Wojtek; Go1's linear ranges are scaled
by $s=2.5$. The control step is 0.02\,s and an episode on a level lasts
300 steps (6\,s). The actor is blind (no terrain perception); the
critic additionally sees a clean $5\times5$ height scan ahead of the robot
(asymmetric actor--critic), so value-based level scores can distinguish
terrains. Training applies physics randomisation~\citep{tobin2017domain}
per environment (Wojtek: terrain friction, base and link masses,
actuator gain and damping scale; Go1: Playground's randomisation).

\paragraph{Arena.}
Levels are rendered into a square heightfield around a flat spawn disc.
In MJX the heightfield data is an array leaf of the model, so each
parallel environment can be reset onto a different level without
recompilation, which makes per-environment level swapping feasible at RL
scale (details in Appendix~\ref{app:render}).

\paragraph{Scaling.}
We set $s = 2.5$, approximately the ratio of Go1's standing height to
Wojtek's ($0.258/0.10$).
Every length in the benchmark (arena, terrain features, gap depth and
period, commands, height-scan grid and success distance) is multiplied by
$s$ for Go1; slope angles are unchanged. Capability envelopes are reported
in Wojtek units so the two robots can be compared directly.

\subsection{Level space}
\label{sec:levelspace}

A level is a vector $\theta\in[0,1]^8$ plus a noise seed. Each coordinate
maps linearly to a physical range (Table~\ref{tab:levelspace}). Features are
measured radially from the edge of the spawn disc, or are seeded random
fields, so the level does not depend on the robot's heading. The rendered
height is the sum of a radial slope, radial stairs, smooth noise, and
piecewise-constant blocks, with gap trenches (0.15\,m deep, every 0.6\,m)
cut in. Each parameter has a \emph{flat} value (zero slope, zero step,
zero amplitude, zero gap, zero block height); because slope and step ranges
are asymmetric, flat is not the centre of the cube.

\begin{table}[t]
  \centering\small
  \caption{Level space (v10), Wojtek units; Go1 multiplies every length by
  $s=2.5$. A level is $\theta\in[0,1]^8$ plus a noise seed.}
  \label{tab:levelspace}
  \begin{tabular}{lll}
    \toprule
    Parameter & Range & Pilot (v8) range \\
    \midrule
    Noise amplitude & 0--0.10\,m & same \\
    Noise wavelength & 0.2--1.0\,m & same \\
    Radial slope & $-20^\circ$ to $+18^\circ$ & $\pm20^\circ$ \\
    Stair height & $-0.08$ to $+0.05$\,m & $\pm0.08$\,m \\
    Stair width & 0.15--0.60\,m & same \\
    Gap width & 0--0.06\,m & 0--0.25\,m \\
    Block height & 0--0.06\,m & 0--0.10\,m \\
    Block size & 0.15--0.60\,m & same \\
    \bottomrule
  \end{tabular}
\end{table}

\paragraph{Calibration.}
A level space is only useful for comparing curricula if uniform sampling is
neither trivially easy (nothing to curate) nor mostly unsolvable (curation
and DR both train on impossible levels). Our pilot space (v8) was the
latter: the best pilot policies solved about 8\% of uniform levels, and
about 92\% of uniform levels were unsolved by every method. We recalibrate
with the pilot policies themselves: for each feature we measure the best
v8 policies' capability envelope, the largest feature value still passed
with 50\% success when that feature alone is active
(\S\ref{sec:protocol}), and cap the uphill slope, step-up height, gap width
and block height at that envelope. Estimated from the subset of v8
evaluation levels that fall inside the new box, v8 policies solve about
39\% of uniform v10 levels. Because features compound, uniform levels still
range from trivial to unsolved. We regard this procedure, sizing the space
to the capability envelopes of a reference policy, as part of the protocol:
it is cheap, uses only the evaluation tooling, and transfers to any
parametric level space.
\todo{Report the v10 uniform success of DR as a check that the calibration
target was met, once results land.}

\section{Student and Methods}
\label{sec:methods}

\subsection{Student and training loop}
All methods train the same student: brax's PPO
networks, loss and hyperparameters from MuJoCo Playground's locomotion
recipe (policy and value MLPs $512\text{-}256\text{-}128$ with ELU, running
observation normalisation, 20-step unrolls, 4 epochs $\times$ 32
minibatches, $\gamma=0.97$, $\lambda=0.95$, learning rate $3\times10^{-4}$,
entropy cost $10^{-2}$, clipping 0.2). Brax's own training loop fixes each
environment's model for a whole run, so we drive the same components from a
custom loop. One outer \emph{iteration} runs 256 lanes $\times$ 32
environments; every lane holds one level, all 32 of its environments are
reset onto it with a fresh physics-randomisation draw, and the batch is
stepped for 300 steps in 20-step unrolls, with a PPO update after each
unroll when the iteration trains. An iteration is therefore
$256\cdot32\cdot300\approx2.46$M environment steps. Each lane's score is
either the positive value loss, the mean positive GAE advantage over the
iteration~\citep{jiang2021plr}, or the learnability $p(1-p)$, where $p$ is
the fraction of the lane's environments that never fell and travelled at
least 1\,m ($\times s$) from spawn~\citep{rutherford2024noregrets}.
Runs last 300M environment steps for Wojtek and 1B for Go1, with three
seeds per method and robot.

\subsection{Methods}
\label{sec:methodlist}

\paragraph{DR.} Every iteration trains on 256 fresh uniform levels.

\paragraph{DR-easy.} DR on ACCEL's easy generator: a uniform genome shrunk
toward the flat genome by a factor 0.1. It isolates the effect of
``training on easy levels'' from the effect of curation.

\paragraph{\plr~\citep{jiang2021replay}.} A buffer of the 4000
highest-scoring levels (positive value loss). Once it holds at least 256
levels, an iteration replays with probability 0.8: it samples 256 levels by
rank prioritisation (temperature 0.3) mixed with staleness (weight 0.3),
trains on them and refreshes their scores. Otherwise it scores 256 fresh
uniform levels without training and adds them to the buffer.

\paragraph{ACCEL~\citep{parkerholder2022accel}.} \plr\ with two changes:
new levels come from the easy generator instead of the uniform one, and
every replayed level is mutated by perturbing two randomly chosen
parameters by $\mathcal{U}(-0.1,0.1)$ (clipped to $[0,1]$, seed kept). The
next non-training iteration scores the children of the most recent replay
batch (children of earlier consecutive replay batches are discarded), and
fresh easy levels fill the remaining lanes.

\paragraph{ACCEL-SFL.} ACCEL with two changes: the learnability score
$p(1-p)$ in place of positive value loss, and a ``grow'' mutation that,
after the local edit, scales the genome's distance from flat by
$\mathcal{U}(1,1.3)$. Unlike SFL~\citep{rutherford2024noregrets}, which
samples uniformly and keeps a learnability buffer, it keeps ACCEL's
generator and replay. ACCEL and ACCEL-SFL therefore differ in two
switches, score and mutation, and a difference between them cannot be
attributed to the score alone. The pilot (\S\ref{sec:pilot}) ran the
missing arm, ACCEL-grow (grow mutation, positive value loss); the main
study does not.

\paragraph{PAIRED~\citep{dennis2020paired}.} A learned generator maps a
16-dimensional Gaussian latent through an MLP to a sigmoid-squashed
Gaussian over the 8 parameters; a random seed is appended. Each iteration
it proposes 256 levels, on which a protagonist and an antagonist (two
independent brax students) both train. The generator's reward per level is
$\max(0,\ \max R_{\text{ant}} - \overline{R}_{\text{pro}})$, the best
antagonist return minus the mean protagonist return over the lane, and it
is updated by single-step PPO (5 epochs, entropy coefficient 0.1). We do
not use the behaviour-cloning distillation between the students
proposed by~\citet{mediratta2023stabilizing}. The budget
counts protagonist steps only, so PAIRED uses twice the simulation of the
other methods.

\paragraph{Budget accounting.}
All methods get the same number of environment steps. Non-training
iterations of \plr, ACCEL and ACCEL-SFL count toward the budget, so these
methods take fewer gradient steps than DR. We report wall-clock and
simulated steps per method in Table~\ref{tab:compute}.

\section{Evaluation Protocol}
\label{sec:protocol}

Each final policy is evaluated deterministically, without physics
randomisation and without early episode cuts. An episode \emph{succeeds}
if the robot never falls and ends at least 1\,m ($\times s$) from spawn
after 300 steps (6\,s). Every level is rolled out in two independent
passes of 4 episodes each, giving a per-level success rate per pass.

\paragraph{Uniform levels.} 2048 levels drawn uniformly from the level
space, the training distribution of DR. We report mean success with a 95\%
bootstrap confidence interval over levels.

\paragraph{Cross-fitted CVaR.} Robustness is the mean success on the worst
$\alpha$ fraction of levels ($\alpha\in\{0.1, 0.25\}$), in the spirit of
the adversarial evaluation of~\citet{rutherford2024noregrets}. Ranking and
scoring levels on the same episodes biases this estimate low through
selection noise, so we rank levels on one pass and score them on the other,
in both directions, and average. In a simulation with 2000 levels of true
success 0.5 and 4 episodes per level, same-pass \cvar{0.25} reads 0.19
while the cross-fitted estimate reads 0.49.

\paragraph{Level reliability.} The Pearson correlation of per-level success
between the two passes. A high value means difficulty is a property of the
level rather than of the episode, which is a precondition for any curation
method to have something to select. On our earlier control-point
heightfield space this correlation was only 0.02--0.17; on the parametric
pilot space it was 0.41--0.74 (Appendix~\ref{app:pilot}).

\paragraph{Difficulty sweep.} The same 2048 genomes shrunk toward flat by
$d\in\{0.25,0.5,0.75,1.0\}$. We report $d=0.5$ and $d=1.0$; $d=0.25$ is
solved by almost any policy trained near flat terrain and does not separate
methods. A flat-terrain reference ($d=0$) checks that a method has not
traded basic walking for terrain skills.
\todo{$d=0$ is not in \texttt{eval\_brax.py} yet; add it or drop this
sentence.}

\paragraph{Solvable subset.} Mean success on the uniform levels that at
least one evaluated policy solves, which removes levels no method can pass
from the denominator.
\todo{Computed offline from the per-level \texttt{.npz} outputs; define
``solves'' (any episode vs.\ $\geq50\%$) before reporting.}

\paragraph{Capability envelopes.} For each of seven features (noise,
slope up, slope down, stairs up, stairs down, gaps, blocks), the feature
ramps from off to the end of its range in six steps with all other
features off, 64 seeded levels per step. The \emph{envelope} is the
largest feature magnitude still passed with 50\% success (linearly
interpolated), the \emph{AUC} the mean success over the ramp.
Envelopes are in Wojtek units for both robots.

\paragraph{Held-out suite.} Seven structured terrain families no method
trains on, each at four strengths $k=1..4$ and 64 instances, radially
symmetric like the training levels (Wojtek units, $\times s$ for Go1):
stairs up and down ($2k$\,cm steps every 0.3\,m), slopes up and down
($5k^\circ$), trenches (10\,cm deep, $5k$\,cm wide, every 0.5\,m), waves
($\pm2k$\,cm, 0.6\,m wavelength) and random blocks ($\pm2k$\,cm, 0.3\,m
cells). We report mean success over the 28 cells and the CVaR over cells.

\paragraph{Sim2sim transfer.} Policies trained in MJX are re-evaluated in
CPU MuJoCo under five physics variants: \texttt{mjx\_match} (settings
matching training, a sanity check), \texttt{newton\_elliptic},
\texttt{half\_dt}, \texttt{full\_collision} and \texttt{soft\_contact}.
\todo{Describe each variant (solver, friction cone, timestep, collision
set, contact softness) from \texttt{src/scripts/eval\_sim2sim.py} once it
exists, and state which level sets it runs.}

\paragraph{Aggregation.} Per-seed metrics are averaged over three seeds;
per-level outputs are stored so that interquartile means and
stratified bootstrap intervals across seeds~\citep{agarwal2021precipice}
can be computed.
\todo{Decide between seed mean $\pm$ std and IQM with stratified bootstrap
for the main table; three seeds is few for either.}

\section{Results}
\label{sec:results}

\todo{v10 results are not available. Everything below the pilot summary is
the planned analysis: hypotheses carried over from the pilot, not
outcomes. Replace each hypothesis paragraph with the measured result and
delete the ones the data does not address.}

\subsection{Pilot study}
\label{sec:pilot}
We first ran all methods on Wojtek over a wider level space (v8,
Table~\ref{tab:levelspace}, right column), with two additional ACCEL
variants and a DR variant with a random per-level difficulty (DR-mixed);
full numbers are in Appendix~\ref{app:pilot}. Four observations shaped the
benchmark:
\begin{itemize}
  \item \textbf{The starting distribution dominated.} Methods that start
    from easy levels (ACCEL, its variants and DR-easy) solved 7.7--8.2\% of
    uniform levels; the others (DR-mixed, \plr, PAIRED, DR) solved
    0.4--3.5\%.
  \item \textbf{Curation rescued a bad generator.} \plr\ and DR see the
    same uniform levels, but \plr\ solved 2.2\% vs.\ 0.4\% of uniform
    levels and 9.8\% vs.\ 1.0\% at $d=0.5$.
  \item \textbf{ACCEL matched DR-easy.} 7.7\% vs.\ 8.0\% uniform success,
    \cvar{0.25} 2.1\% vs.\ 2.2\%. ACCEL-SFL (learnability score plus grow
    mutation) gave the only gain over DR-easy, and only in the tail
    (\cvar{0.25} 2.8\%; mean 8.2\% vs.\ 8.0\%). ACCEL and
    \plr\ differ in both generator and mutation, so the pilot does not
    isolate the effect of mutation.
  \item \textbf{The space was too hard.} About 92\% of uniform levels were
    solved by no method, which compresses differences between the best
    methods and means DR trains mostly on levels it cannot learn from. This
    motivated the calibration of \S\ref{sec:levelspace}.
\end{itemize}

\subsection{Main comparison}
\label{sec:main}

\begin{table}[t]
  \centering\small
  \caption{Final evaluation on the v10 level space, both robots, mean over
  3 seeds. Uniform success with 95\% bootstrap CI, cross-fitted
  \cvar{0.1} and \cvar{0.25}, success at $d=0.5$, level reliability and
  held-out suite success.}
  \label{tab:main}
  \resulttable{main}{Main results table (\texttt{tables/main.tex}):
  rows = 6 methods $\times$ 2 robots; columns = uniform success [CI],
  \cvar{0.1}, \cvar{0.25}, $d=0.5$, $d=1.0$, level corr., suite success,
  solvable-subset success.}
\end{table}

\begin{figure}[t]
  \centering
  \placeholder{Learning curves: deterministic success on 512 fixed eval
  levels vs.\ environment steps, one panel per robot (and per
  $d\in\{0.5,1.0\}$), one line per method, mean over 3 seeds with min--max
  band. Source: \texttt{eval/d<d>/success\_rate} logged every 10
  iterations.}
  \caption{Learning curves.}
  \label{fig:curves}
\end{figure}

Hypotheses carried over from the pilot, to be tested on v10:
\begin{description}
  \item[H1 (generator).] Methods with an easy generator (DR-easy, ACCEL,
    ACCEL-SFL) outperform those with a uniform one (DR, \plr, PAIRED) on
    uniform success. With a calibrated space DR should no longer be near
    zero, so the gap is expected to shrink.
    \todo{result}
  \item[H2 (curation).] \plr\ improves over DR at the same generator.
    \todo{result; report the ratio with CIs on both robots}
  \item[H3 (mutation).] ACCEL improves over DR-easy, the only comparison
    that holds the generator fixed and adds curation and mutation.
    \todo{result}
  \item[H4 (score).] Relative to ACCEL, ACCEL-SFL gains more in the
    tail (\cvar{0.1} and \cvar{0.25}) than in mean success.
    \todo{result; confounded with the grow mutation
    (\S\ref{sec:methodlist})}
  \item[H5 (PAIRED).] PAIRED without distillation is the weakest UED
    method despite twice the simulation.
    \todo{result}
  \item[H6 (morphology).] The method ranking on Go1 matches Wojtek's.
    \todo{result; rank correlation across robots}
\end{description}

\subsection{Sim2sim transfer}
\label{sec:sim2sim}
\begin{table}[t]
  \centering\small
  \caption{Sim2sim transfer: success of MJX-trained policies in CPU
  MuJoCo under each physics variant, mean over 3 seeds.}
  \label{tab:sim2sim}
  \resulttable{sim2sim}{Sim2sim table (\texttt{tables/sim2sim.tex}):
  rows = methods $\times$ robots; columns = MJX reference,
  \texttt{mjx\_match}, \texttt{newton\_elliptic}, \texttt{half\_dt},
  \texttt{full\_collision}, \texttt{soft\_contact}.}
\end{table}
\begin{description}
  \item[H7 (solver sensitivity).] The closed-chain Wojtek loses more
    success than the open-chain Go1 when the solver or contact model
    changes, since its legs depend on the equality-constraint closure.
    \todo{result}
  \item[H8 (curriculum).] Curricula that train on harder terrain transfer
    better across physics variants.
    \todo{result; define ``harder'' via the composition logs of
    Figure~\ref{fig:composition}}
\end{description}

\subsection{What the curricula train on}
\begin{figure}[t]
  \centering
  \placeholder{Curriculum composition over training: per method, the mean
  physical value of each of the 8 level parameters among training levels
  (\texttt{level/<param>}) vs.\ environment steps, plus the genome distance
  from flat. One small multiple per parameter, one line per method; Wojtek
  top row, Go1 bottom row.}
  \caption{Curriculum composition.}
  \label{fig:composition}
\end{figure}
\todo{Describe which features each curriculum grows first (e.g.\ whether
ACCEL's buffer moves away from the easy generator at all, which it did not
on the earlier control-point space) once \texttt{level/<param>} is
logged.}

\subsection{Capability envelopes}
\begin{table}[t]
  \centering\small
  \caption{Capability envelopes: largest feature magnitude passed with
  50\% success, Wojtek units, mean over seeds.}
  \label{tab:envelopes}
  \resulttable{envelopes}{Envelope table (\texttt{tables/envelopes.tex}):
  rows = methods $\times$ robots; columns = rough (m), slope up/down
  ($^\circ$), stairs up/down (m), gaps (m), blocks (m).}
\end{table}

Per-feature envelopes are plotted in Figure~\ref{fig:envelopes}
(Appendix~\ref{app:extra}).
\todo{Which features separate methods; whether Go1 and Wojtek fail on the
same features (e.g.\ step-up for the small closed-chain robot).}

\section{Discussion and Limitations}
\label{sec:discussion}

\todo{Write after v10 results. Planned points, kept only if the data
supports them:}
\begin{itemize}
  \item \textbf{Generator vs.\ curation.} Whether the pilot's ordering
    (easy generator $\gg$ curation $>$ none) survives calibration, i.e.\
    whether UED's benefit on 3D locomotion is mostly ``start easy''.
  \item \textbf{Regret estimates on locomotion.} Positive value loss
    depends on the critic, which here sees a privileged height scan;
    whether learnability is a better score for a continuous-control
    student with a binary success criterion.
  \item \textbf{Closed vs.\ open chain.} Whether the small closed-chain
    robot and Go1 favour the same methods and fail on the same features.
\end{itemize}

\paragraph{Limitations.}
\begin{itemize}
  \item Three seeds per method and robot; small differences (such as the
    pilot's tail gain of ACCEL-SFL) need more seeds to confirm.
  \item One student (brax PPO) and one budget per robot; methods that
    take fewer gradient steps (\plr, ACCEL) or simulate twice (PAIRED)
    are compared at equal environment steps, not equal compute.
  \item The level space is parametric and radially symmetric; it does not
    cover unstructured terrain maps or heading-dependent obstacles.
  \item The calibration uses pilot policies from the same benchmark, so
    the level space is tuned to this student and these robots.
  \item Simulation only: transfer is tested sim-to-sim (MJX to CPU
    MuJoCo), not on the physical Wojtek.
  \item PAIRED is run without the behaviour-cloning distillation
    of~\citet{mediratta2023stabilizing}, and ACCEL-SFL changes two
    switches relative to ACCEL.
\end{itemize}

\paragraph{Conclusion.}
\todo{Two sentences once the headline result is known.}


\bibliographystyle{plainnat}
\bibliography{refs}

\clearpage
\appendix
\section{Pilot Study (v8 Level Space, Wojtek)}
\label{app:pilot}

The pilot used the level space of Table~\ref{tab:levelspace} (right
column), Wojtek only, 300M environment steps, 3 seeds, and the evaluation
protocol of \S\ref{sec:protocol}. Besides the six methods of
\S\ref{sec:methods} it included DR-mixed (uniform genome shrunk toward flat
by a per-level difficulty $d\sim\mathcal{U}(0,1)$) and ACCEL-grow (ACCEL
with the grow mutation, positive value loss score); ACCEL-SFL is
ACCEL-grow with the learnability score.

\begin{table}[ht]
  \centering\small
  \caption{Pilot (v8) final evaluation, mean over 3 seeds. Success on
  uniform levels, cross-fitted \cvar{0.25}, pass-to-pass level success
  correlation, success at difficulty $d=0.5$ and $d=1.0$, and mean success
  on the held-out suite.}
  \label{tab:pilot}
  \begin{tabular}{lcccccc}
    \toprule
    Method & Uniform & \cvar{0.25} & Level corr. & $d=0.5$ & $d=1.0$ & Suite \\
    \midrule
    ACCEL-SFL  & .082 & .028 & .667 & .337 & .077 & .521 \\
    DR-easy    & .080 & .022 & .716 & .334 & .073 & .548 \\
    ACCEL-grow & .078 & .020 & .716 & .329 & .073 & .556 \\
    ACCEL      & .077 & .021 & .737 & .332 & .070 & .559 \\
    DR-mixed   & .035 & .012 & .689 & .181 & .030 & .445 \\
    \plr       & .022 & .009 & .692 & .098 & .018 & .382 \\
    PAIRED     & .016 & .006 & .673 & .076 & .013 & .320 \\
    DR         & .004 & .002 & .414 & .010 & .002 & .120 \\
    \bottomrule
  \end{tabular}
\end{table}

\todo{Add seed standard deviations and the v8 capability envelopes
(\texttt{paper\_tables.py --prefix wojtek=v8\_param}).}

\section{Arena and Level Rendering}
\label{app:render}
A level is rendered into a square heightfield of half-width 3.0\,m at
0.05\,m resolution ($121\times121$ nodes) with heights clipped to
$\pm0.3$\,m (all $\times s$ for Go1). A flat spawn disc of radius 0.55\,m
blends into the terrain over 0.2\,m. In MJX the heightfield data is an
array leaf of the model, so each of the parallel environments can be reset
onto a different level without recompilation; this is what makes
per-environment level swapping feasible at RL scale. We run MJX on its JAX
backend, because the Warp backend (at the pinned version) collides every
environment in a batch with the first environment's heightfield.

Let $r$ be the distance from the arena centre, $r_0$ the spawn-pad radius,
$w$ the taper width, $\delta=\max(r-r_0-w,\,0)$ the radial distance past
the pad and $\tau(r)=\mathrm{clip}((r-r_0)/w,0,1)$. With the physical
parameters of Table~\ref{tab:levelspace},
\begin{equation*}
  h(x,y) = \tau(r)\Big[\tan(\phi)\,\delta
    + h_s\Big\lceil\frac{\delta}{w_s}\Big\rceil
    + A\,n_\lambda(x,y) + h_b\,b_c(x,y)
    - D\,\mathbf{1}\big[\delta>0,\ (\delta \bmod P) < g\big]\Big],
\end{equation*}
where $\phi$ is the slope, $h_s,w_s$ stair height and width, $A,\lambda$
noise amplitude and wavelength, $n_\lambda$ a seeded bilinear noise field
in $[-1,1]$ with control points every $\lambda$, $b_c$ a seeded
piecewise-constant field in $[-1,1]$ with cells of size $c$ and height
$h_b$, and $g$ the gap width with depth $D=0.15$\,m and period $P=0.6$\,m.
Heights are clipped to the arena amplitude. For Go1 every length is
multiplied by $s=2.5$.

\section{Additional Results}
\label{app:extra}

\begin{figure}[t]
  \centering
  \placeholder{Capability envelopes per feature: grouped bars (or one radar
  per robot), envelope normalised by the feature's range, one bar per
  method, seed min--max as error bars.}
  \caption{Capability envelopes.}
  \label{fig:envelopes}
\end{figure}

\section{Hyperparameters and Compute}
\label{app:compute}

\begin{table}[ht]
  \centering\small
  \caption{Training hyperparameters shared by all methods.}
  \label{tab:hparams}
  \begin{tabular}{ll}
    \toprule
    Lanes $\times$ envs per lane & $256\times32$ \\
    Steps per level (episode) & 300 (6\,s) \\
    Env steps & 300M (Wojtek), 1B (Go1) \\
    PPO & brax; 20-step unroll, 4 epochs $\times$ 32 minibatches \\
    $\gamma$, $\lambda$, lr, entropy, clip & 0.97, 0.95, $3\times10^{-4}$, 0.01, 0.2 \\
    Networks & MLP 512-256-128, ELU, obs.\ normalisation \\
    Buffer (\plr, ACCEL) & 4000 levels; replay prob.\ 0.8 \\
    Replay sampling & rank (temp.\ 0.3) + staleness (0.3) \\
    Easy generator & shrink toward flat $\times0.1$ \\
    Mutation (ACCEL) & 2 parameters, $\pm0.1$ \\
    Grow (ACCEL-SFL) & distance from flat $\times\,\mathcal{U}(1,1.3)$ \\
    PAIRED generator & latent 16, hidden 128, lr $10^{-4}$, entropy 0.1 \\
    \bottomrule
  \end{tabular}
\end{table}

\begin{table}[ht]
  \centering\small
  \caption{Training compute per run, mean $\pm$ std over seeds. GPU-h: wall
  clock on one GPU; sim steps count every simulated environment step
  (PAIRED's antagonist, the non-training iterations of \plr\ and ACCEL);
  student steps are those that fed a student PPO update.}
  \label{tab:compute}
  \resulttable{compute}{Compute table (\texttt{tables/compute.tex}):
  rows = robots $\times$ methods; columns = GPU-h, sim steps, student
  steps.}
\end{table}
\todo{State the hardware (WCSS Lem, one H100 per run) once the compute
table confirms it for v10.}


\end{document}
